\chapter{Avaliar Pedidos e Investigar Falhas}
\label{ch:avaliacao}

Imagina duas respostas a um pedido de reembolso. A primeira cita a política e promete 48 horas. A segunda cita a mesma política e encaminha a reclamação. Se a política deixa o prazo por definir, a primeira resposta falha. Mais adiante, veremos este erro num ensaio real, com a resposta guardada.

Para avaliar, começa pela decisão que a resposta pode afetar. Define o que seria uma saída aceitável e como descobririas uma falha. Depois fixa as entradas e compara versões. Uma troca de prompt, de modelo ou de fonte pode mudar a resposta. O teste precisa de saber que peça mudou.

\section{Definir o que conta como resposta aceitável}

Começa pelo risco da decisão. Uma resposta pode cumprir o formato, usar um tom adequado e prometer um prazo que a fonte não autoriza. Regista o erro que reprova o caso, a evidência para o detetar e as preferências que ainda permitem aceitar a resposta.

\begin{promptlabbox}{Definir critérios para uma tarefa tua}
\textbf{Objetivo e contexto:} escolher condições de aceitação antes de calcular métricas. Usa uma tarefa delimitada, como resumir uma reunião ou classificar mensagens.

\textbf{Procedimento:} escreve um critério factual, um de formato e um de ação. Para cada um, cria um caso que passe e outro que falhe. Indica se a falha obriga a parar ou se pode ser revista.

\textbf{Resultado esperado e validação:} outra pessoa consegue aplicar os critérios e indicar a fonte ou o trecho que sustenta cada decisão.

\textbf{Falhas possíveis:} tratar preferência de estilo como erro crítico, usar um critério que nenhum caso pode contrariar ou escolher a medida depois de ver os resultados.
\end{promptlabbox}

\section{Preparar um conjunto com respostas de referência}

Um \emph{golden set} reúne entradas e resultados esperados que podemos conferir. Inclui casos correntes, ambiguidades, informação em falta e erros com consequências. Se a tarefa é classificar, guarda a categoria e o motivo. Se a resposta tem vários campos, define o que deve acontecer em cada um. Estes dois registos fictícios mostram formas de guardar os casos:

\begin{lstlisting}
{
  "id": "exemplo_01",
  "input": "Quero confirmar o estado de LAV-001",
  "expected_label": "estado_reclamacao"
}
\end{lstlisting}

\begin{lstlisting}
{
  "id": "exemplo_07",
  "input": "Fui cobrado duas vezes",
  "expected_output": {
    "tipo_problema": "pagamento",
    "urgencia": "alta",
    "precisa_resposta_humana": true
  }
}
\end{lstlisting}

No segundo registo, «alta» e o encaminhamento pressupõem uma regra de triagem que tens de definir. A categoria pode ser confirmada pelas palavras da mensagem. A urgência exige o critério da tarefa. Guarda essa diferença no resultado esperado e pede revisão a quem conhece o processo.

\section{Escolher uma medida que observe a falha}

Uma etiqueta pode ser comparada diretamente. Numa extração, podes verificar campos e tipos. Numa resposta livre, várias formulações podem estar certas. Nesse caso, mede a cobertura das ideias necessárias e confirma se cada afirmação tem suporte. A proximidade entre textos pode ajudar a escolher o que ler primeiro, mas duas frases parecidas podem discordar num «não» ou num número.

\section{Medir acerto, cobertura e custo do erro}
\label{sec:metricas}
Imagina cinco pedidos que exigem atenção. O sistema sinaliza três desses cinco e ainda um pedido normal. Das quatro sinalizações, três estavam certas. A \emph{precisão} é $3/4=75\%$. Dos cinco pedidos que exigiam atenção, três foram sinalizados. A \emph{sensibilidade} é $3/5=60\%$. O numerador é o mesmo, mas a primeira pergunta conta as sinalizações e a segunda conta os casos reais.

\textbf{Aprofundamento opcional.} Nas fórmulas, $TP$ conta os verdadeiros positivos (3 no exemplo), $FP$ os falsos positivos (1) e $FN$ os falsos negativos (2). A precisão é $TP/(TP+FP)$ e a sensibilidade é $TP/(TP+FN)$. Regista como o cálculo trata denominadores iguais a zero. A medida F1 combina precisão $p$ e sensibilidade $s$ como $2ps/(p+s)$ quando esta expressão está definida. Com várias classes, diz se calculas por classe e fazes a média (\emph{macro}) ou se juntas as contagens antes de calcular (\emph{micro}).

Na Lavandaria, encaminhar todos os pedidos para uma pessoa evita algumas decisões automáticas erradas, mas não significa resolver a tarefa. Mede cobertura automática (respostas emitidas/pedidos), acerto entre respostas emitidas e erros graves por pedido. Define antecipadamente se uma abstinência correta conta como sucesso operacional e apresenta-a separadamente. Num exemplo sintético de dez pedidos, A responde a todos e acerta oito. B responde a cinco e acerta os cinco. B tem maior acerto entre respostas, 100\% contra 80\%, e menor cobertura, 50\% contra 100\%. Para escolher, falta saber o custo dos erros e dos cinco encaminhamentos.

Dez respostas, oito certas, dão uma proporção observada de $0.8$. Não dão «80\% de fiabilidade» para todos os clientes. Pelo intervalo de Wilson a 95\%, sob um modelo binomial com casos independentes e representativos, o intervalo é aproximadamente $[0.490,0.943]$~\citep{wilson1927}. Dez variações quase iguais não justificam a mesma generalização. A demonstração reproduz o cálculo em Python.

\section{Comparar versões e preservar o teste}

Numa regressão, aplicas a versão antiga e a nova aos mesmos casos e verificas o que melhorou, o que se manteve e onde surgiu uma falha. Identifica o modelo, o pedido, a configuração e as fontes. Usa o ciclo do capítulo~\ref{ch:mentalidade} para desenvolver. Guarda casos diferentes para avaliar a versão fechada. Quando uma falha reservada orienta uma correção, o caso pode ficar como regressão e a próxima estimativa exige casos novos.

Se a pergunta é sobre o efeito de uma mudança no serviço real, uma comparação A/B pode atribuir as versões aleatoriamente. Escolhe antes a unidade de atribuição, por exemplo a conta, e a medida principal. Se a mesma pessoa alternar entre versões, uma interação pode influenciar a seguinte. Regista também falhas, encaminhamentos e tempo até à resolução. Um lançamento canário tem outra pergunta: com exposição limitada, há um sinal que obriga a travar? O capítulo~\ref{ch:llmops} distingue as duas decisões.

\begin{figure}[htbp]
\centering
\begin{tikzpicture}[node distance=6mm and 8mm, every node/.style={draw,align=center,text width=39mm,minimum height=10mm,font=\small},>=Stealth]
\node (dev) {Desenvolvimento e critérios};
\node (freeze) [below=of dev] {Versão fechada};
\node (test) [below left=of freeze] {Teste reservado e repetições};
\node (safe) [below right=of freeze] {Segurança e casos críticos};
\node (gate) [below=27mm of freeze] {Decisão documentada};
\draw[->](dev)--(freeze);\draw[->](freeze)--(test);\draw[->](freeze)--(safe);\draw[->](test)--(gate);\draw[->](safe)--(gate);
\end{tikzpicture}
\caption{Avaliação antes da passagem a produção. O teste reservado não é um catálogo de exemplos para copiar para o prompt.}
\label{fig:avaliacao}
\end{figure}

\section{Investigar uma falha antes de mudar o pedido}

Guarda a entrada, o resultado esperado e o ponto onde divergiu. Uma data errada pode vir de um documento desatualizado, de um cálculo, de uma interpretação ou da redação. Para localizar o problema, altera uma peça e repete o caso. Liga a alteração ao resultado e a uma versão que possas recuperar.

\begin{promptlabbox}{Investigar falhas que parecem iguais}
\textbf{Objetivo e contexto:} distinguir erros com a mesma aparência e causas diferentes. Usa falhas registadas ou constrói três casos fictícios.

\textbf{Procedimento:} escolhe um prazo antigo recuperado da fonte errada, um prazo mal lido e um prazo acrescentado sem fonte. Para cada caso, indica a evidência que permitiria separar as hipóteses. Corrige uma causa e volta a correr os casos com a mesma regra.

\textbf{Resultado esperado e validação:} a correção muda o resultado esperado no caso afetado e conserva os resultados certos nos restantes. Guarda o que observaste e o que continua por explicar.

\textbf{Falhas possíveis:} aumentar o pedido sem localizar a causa, avaliar só o caso corrigido ou perder a versão que produziu a falha.
\end{promptlabbox}

\section{Juízes também precisam de auditoria}
Estudos de avaliação com LLMs documentam efeitos como preferência pela posição de uma resposta e pela verbosidade~\citep{zheng2023}. Baralha a ordem, oculta a identidade das versões, usa critérios ancorados em exemplos e compara uma amostra com avaliação humana. Uma resposta mais longa não é automaticamente mais correta. A nota do juiz é uma medição imperfeita, não uma autoridade que dispensa revisão.

\section{Avaliar o sistema que mudou, não apenas a resposta}
Numa skill, separa descoberta e ativação, cumprimento do procedimento, execução dos scripts e qualidade do resultado. Num sistema multiagente, acrescenta comunicação, fusão, orçamento, falhas parciais e recuperação. Um teste com adaptador determinístico verifica o controlador; uma avaliação real com LLM precisa dos pedidos e respostas desse modelo.

Compara também versões em que retiras uma componente: pesquisa, revisão ou especialista. Esta ablação ajuda a perceber se o ganho depende dessa peça ou de contexto e recursos adicionais. Conserva condições e custos visíveis. Um conjunto público repetidamente usado em desenvolvimento pode deixar de estimar comportamento em situações novas; não basta chamar-lhe benchmark para o tornar independente.

\section{Um ensaio real, incluindo o que correu mal}
\label{sec:ensaio-real}
Em 3 de novembro de 2026 foram executadas 52 inferências locais com o modelo oficial Qwen2.5-0.5B-Instruct, quantização Q4\_K\_M, e llama.cpp build b11960~\citep{qwengguf,llamacpp}. O repositório conserva os pedidos, respostas brutas, revisões, hashes, configurações e avaliação em \texttt{experiencias/llm\_local}. Os casos e soluções foram definidos antes das chamadas; não se voltou a afinar o prompt para apagar as falhas.

São três comparações pequenas, cada uma com duas seeds, temperatura 0.2 e contexto novo por chamada. A classificação usa seis mensagens; a resposta documental, três pedidos; a duração, quatro problemas. Na última comparação, ambas as condições usam o mesmo esquema JSON; só muda uma instrução de decomposição, sem pedir uma explicação. Não se reproduz o método few-shot CoT dos estudos originais.

\begin{table}[htbp]
\small
\begin{tabularx}{\textwidth}{@{}lXX@{}}
\toprule
Tarefa e critério & Primeira condição & Segunda condição \\
\midrule
Classificação: etiqueta exata & Zero-shot: 6/12 & Few-shot: 11/12 \\
Documentos: JSON válido & Só prompt: 6/6 & Esquema no runtime: 6/6 \\
Documentos: contrato e excertos & Só prompt: 0/6 & Esquema no runtime: 1/6 \\
Duração: resultado exato & Direta: 1/8 & Decomposição: 0/8 \\
\bottomrule
\end{tabularx}
\caption{Contagens observadas no ensaio local. Cada denominador inclui duas execuções dos mesmos casos; não representa outros tantos casos independentes.}
\label{tab:ensaio-local}
\end{table}

O caso mais instrutivo é uma resposta «48 horas» a um pedido de reembolso. Passou o contrato e citou literalmente POL-02, mas essa fonte diz precisamente que não existe prazo fixo nem reembolso automático. A presença de uma citação verdadeira não faz a resposta decorrer dela. Nas seis saídas com esquema, a estrutura respeitou os campos e tipos; cinco falharam a proveniência verificada e a restante falhou o conteúdo. A revisão de conteúdo registada foi feita durante a edição assistida, sem painel humano independente.

O modelo é pequeno, quantizado e não representa os modelos disponíveis em 2026. A amostra é sintética e curta; as repetições partilham casos e modelo. Estes números demonstram o procedimento e os seus limites. Não provam melhoria geral de few-shot, superioridade de um fornecedor, inutilidade de decomposição ou segurança contra ataques. Mesmo quando todos os pedidos terminam sem erro de execução, a tarefa pode falhar. Para calcular durações, continua a haver uma solução bastante competente: fazer a conta em código.
